\documentclass[10pt,a4paper]{article}
\usepackage[left=2cm,right=2cm,top=2cm,bottom=2cm]{geometry}
\usepackage[utf8]{inputenc}
\usepackage[T1]{fontenc}
\usepackage[spanish,es-tabla]{babel}
\usepackage{array}
\usepackage{booktabs}
\usepackage{caption}
\usepackage[table,dvipsnames]{xcolor}
\usepackage{float}
\usepackage{graphicx}
\usepackage{hyperref}
\usepackage{listings}
\usepackage{longtable}
\usepackage{titlesec}

\definecolor{colorIPN}{rgb}{0.5,0.0,0.13}
\definecolor{colorESCOM}{rgb}{0.0,0.5,1.0}
\definecolor{codeBg}{rgb}{0.96,0.97,0.99}

\graphicspath{{capturas/}}

\hypersetup{
    colorlinks=true,
    linkcolor=colorIPN,
    urlcolor=colorESCOM
}

\titleformat{\section}
  {\normalfont\Large\bfseries\color{colorIPN}}
  {\thesection}{1em}{}
\titleformat{\subsection}
  {\normalfont\large\bfseries\color{colorESCOM}}
  {\thesubsection}{1em}{}

\lstset{
    basicstyle=\small\ttfamily,
    numbers=left,
    numberstyle=\tiny,
    frame=tb,
    tabsize=4,
    columns=fixed,
    showstringspaces=false,
    breaklines=true,
    backgroundcolor=\color{codeBg},
    commentstyle=\color{Violet},
    keywordstyle=\color{colorIPN}\bfseries,
    stringstyle=\color{colorESCOM}
}

\newcommand{\cod}[1]{\texttt{\scriptsize #1}}
\newcommand{\captura}[3]{
\begin{figure}[H]
    \centering
    \includegraphics[width=0.92\textwidth,height=0.56\textheight,keepaspectratio]{#1}
    \caption{#2}
    \label{#3}
\end{figure}
}

\begin{document}

\begin{titlepage}
    \centering
    {\huge\bfseries\color{colorIPN}{Instituto Politécnico Nacional}\par}
    {\Large\bfseries\color{colorESCOM}{Escuela Superior de Cómputo}\par}
    \vspace{1cm}
    {\Large\color{colorIPN}{Desarrollo de Aplicaciones Web Client and Backend Development Frameworks}\par}
    \vspace{1.5cm}
    {\huge\bfseries\color{colorESCOM}{Integración de Swagger/OpenAPI}\par}
    \vspace{0.4cm}
    {\Large\bfseries\color{colorIPN}{Sistema de Gestión de Cuartos y Reservaciones}\par}
    \vspace{2cm}
    {\Large\itshape\color{colorIPN}{Profesor: M. en C. José Asunción Enríquez Zárate}\par}
    \vspace{1.5cm}
    {\Large\itshape\color{colorIPN}{Alumna: Estefany Karina Sánchez Calderón}\par}
    {\Large\itshape\color{colorIPN}{estefanysanchez1495@gmail.com}\par}
    {\Large\itshape\color{colorIPN}{Grupo: 7CM3}\par}
    \vfill
    {\large\color{colorIPN}{\today}\par}
\end{titlepage}

\tableofcontents
\pagebreak

\section{Objetivo de la práctica}

El objetivo de la práctica fue integrar Swagger/OpenAPI en el Back End del Sistema de Gestión de Cuartos y Reservaciones, desarrollado con Spring Boot, para generar documentación automática de la API REST y facilitar la prueba de sus endpoints desde una interfaz interactiva.

La integración se enfocó en documentar los módulos principales solicitados por la actividad:

\begin{itemize}
    \item Habitaciones o cuartos.
    \item Reservaciones.
    \item Autenticación y seguridad.
\end{itemize}

\section{Descripción del sistema}

El sistema permite administrar cuartos de hotel y registrar reservaciones asociadas a usuarios autenticados. El Back End expone servicios REST bajo el prefijo \cod{/api/v1} y utiliza Spring Security con autenticación HTTP Basic para proteger las operaciones según el rol del usuario.

\begin{longtable}{|>{\raggedright\arraybackslash}p{0.22\textwidth}|>{\raggedright\arraybackslash}p{0.58\textwidth}|}
\hline
\rowcolor{cyan!20}
\textbf{Módulo} & \textbf{Responsabilidad}\\
\hline
\texttt{Cuarto} & Crear, consultar, actualizar, cambiar disponibilidad y eliminar habitaciones.\\
\hline
\texttt{Reservacion} & Crear, consultar, actualizar, cancelar y eliminar reservaciones, validando fechas y traslapes.\\
\hline
\texttt{Auth} & Consultar la sesión del usuario autenticado mediante Spring Security.\\
\hline
\texttt{Archivo} & Cargar y descargar archivos asociados al sistema.\\
\hline
\end{longtable}

\section{Configuración de Swagger/OpenAPI}

La documentación automática se habilitó mediante la dependencia \cod{springdoc-openapi-starter-webmvc-ui}. Esta librería genera la especificación OpenAPI en formato JSON y publica Swagger UI para visualizar y probar los endpoints.

\begin{lstlisting}[language=XML,caption={Dependencia usada para integrar Springdoc OpenAPI.}]
<dependency>
    <groupId>org.springdoc</groupId>
    <artifactId>springdoc-openapi-starter-webmvc-ui</artifactId>
    <version>2.8.13</version>
</dependency>
\end{lstlisting}

En el archivo \cod{application.properties} se configuraron las rutas de documentación:

\begin{lstlisting}[language=bash,caption={Rutas de Swagger UI y OpenAPI JSON.}]
springdoc.swagger-ui.path=/swagger-ui.html
springdoc.api-docs.path=/v3/api-docs
springdoc.swagger-ui.operationsSorter=method
springdoc.swagger-ui.tagsSorter=alpha
\end{lstlisting}

La clase \cod{ApiConfig} define los metadatos principales de la documentación:

\begin{lstlisting}[language=Java,caption={Configuración general de OpenAPI.}]
@Configuration
public class ApiConfig {
    @Bean
    public OpenAPI myOpenAPI() {
        return new OpenAPI()
            .info(new Info()
                .title("API para la gestion de reservaciones")
                .version("1.0")
                .description("API REST para la gestion de Cuartos y Reservaciones"));
    }
}
\end{lstlisting}

\captura{01-swagger-ui-general.png}{Interfaz Swagger UI del Sistema de Gestión de Cuartos y Reservaciones.}{fig:swagger-general}

\section{Documentación de endpoints}

Los controladores principales usan anotaciones de OpenAPI mediante \cod{@Tag}, lo que permite agrupar las operaciones por módulo dentro de Swagger UI.

\subsection{Módulo de cuartos}

El módulo \cod{Cuarto} se encuentra bajo la ruta \cod{/api/v1/cuartos}. Sus operaciones permiten administrar habitaciones y consultar disponibilidad.

\begin{longtable}{|>{\raggedright\arraybackslash}p{0.18\textwidth}|>{\raggedright\arraybackslash}p{0.46\textwidth}|>{\raggedright\arraybackslash}p{0.25\textwidth}|}
\hline
\rowcolor{cyan!20}
\textbf{Método} & \textbf{Ruta} & \textbf{Uso}\\
\hline
POST & \texttt{/api/v1/cuartos} & Crear cuarto.\\
\hline
GET & \texttt{/api/v1/cuartos} & Listar cuartos.\\
\hline
GET & \texttt{/api/v1/cuartos/disponibles} & Listar cuartos disponibles.\\
\hline
GET & \texttt{/api/v1/cuartos/\{id\}} & Consultar cuarto por identificador.\\
\hline
PUT & \texttt{/api/v1/cuartos/\{id\}} & Actualizar datos de un cuarto.\\
\hline
PATCH & \texttt{/api/v1/cuartos/\{id\}/disponibilidad} & Cambiar disponibilidad.\\
\hline
DELETE & \texttt{/api/v1/cuartos/\{id\}} & Eliminar cuarto.\\
\hline
\end{longtable}

\captura{02-swagger-cuartos.png}{Endpoints del módulo de cuartos documentados en Swagger UI.}{fig:swagger-cuartos}

\subsection{Módulo de reservaciones}

El módulo \cod{Reservacion} se encuentra bajo la ruta \cod{/api/v1/reservaciones}. Incluye validaciones para evitar fechas inválidas y traslapes de reservaciones activas sobre el mismo cuarto.

\begin{longtable}{|>{\raggedright\arraybackslash}p{0.18\textwidth}|>{\raggedright\arraybackslash}p{0.46\textwidth}|>{\raggedright\arraybackslash}p{0.25\textwidth}|}
\hline
\rowcolor{cyan!20}
\textbf{Método} & \textbf{Ruta} & \textbf{Uso}\\
\hline
POST & \texttt{/api/v1/reservaciones} & Crear reservación.\\
\hline
GET & \texttt{/api/v1/reservaciones} & Listar todas las reservaciones.\\
\hline
GET & \texttt{/api/v1/reservaciones/\{id\}} & Consultar reservación por identificador.\\
\hline
GET & \texttt{/api/v1/reservaciones/usuario/\{idUsuario\}} & Consultar reservaciones por usuario.\\
\hline
GET & \texttt{/api/v1/reservaciones/cuarto/\{idCuarto\}} & Consultar reservaciones por cuarto.\\
\hline
PUT & \texttt{/api/v1/reservaciones/\{id\}} & Actualizar fechas.\\
\hline
PATCH & \texttt{/api/v1/reservaciones/\{id\}/cancelar} & Cancelar reservación.\\
\hline
DELETE & \texttt{/api/v1/reservaciones/\{id\}} & Eliminar reservación.\\
\hline
\end{longtable}

\captura{03-swagger-reservaciones.png}{Endpoints del módulo de reservaciones documentados en Swagger UI.}{fig:swagger-reservaciones}

\section{Seguridad y acceso a la documentación}

La API utiliza Spring Security con HTTP Basic. Los usuarios iniciales se crean desde \cod{DataInitializer}:

\begin{longtable}{|c|c|c|}
\hline
\rowcolor{cyan!20}
\textbf{Usuario} & \textbf{Contraseña} & \textbf{Rol}\\
\hline
\texttt{admin} & \texttt{123456} & \texttt{ROLE\_ADMIN}\\
\hline
\texttt{user} & \texttt{123456} & \texttt{ROLE\_USER}\\
\hline
\end{longtable}

La configuración de seguridad permite consultar Swagger UI y OpenAPI JSON sin autenticación, mientras que los endpoints funcionales quedan protegidos por roles.

\begin{lstlisting}[language=Java,caption={Reglas principales de seguridad.}]
.requestMatchers("/swagger-ui/**", "/v3/api-docs/**", "/swagger-ui.html").permitAll()
.requestMatchers(HttpMethod.GET, "/api/v1/cuartos/**").hasAnyRole("ADMIN", "USER")
.requestMatchers(HttpMethod.POST, "/api/v1/cuartos/**").hasRole("ADMIN")
.requestMatchers(HttpMethod.GET, "/api/v1/reservaciones").hasRole("ADMIN")
.requestMatchers(HttpMethod.POST, "/api/v1/reservaciones/**").hasAnyRole("ADMIN", "USER")
.httpBasic(Customizer.withDefaults());
\end{lstlisting}

\captura{06-security-config.png}{Configuración de seguridad para Swagger, cuartos y reservaciones.}{fig:security-config}

\section{Validación de OpenAPI}

La especificación OpenAPI se puede consultar desde:

\begin{center}
\url{http://localhost:8081/v3/api-docs}
\end{center}

Este documento JSON contiene la versión de OpenAPI, los datos generales de la API, las rutas disponibles, los esquemas de entrada y salida, los parámetros y los códigos de respuesta.

\captura{05-openapi-json.png}{Especificación OpenAPI JSON generada automáticamente.}{fig:openapi-json}

\section{Pruebas de funcionamiento}

Se compiló el proyecto con Maven Wrapper usando el comando:

\begin{lstlisting}[language=bash,caption={Compilación del proyecto.}]
./mvnw -DskipTests package
\end{lstlisting}

El resultado fue \cod{BUILD SUCCESS}, confirmando que la integración de Swagger/OpenAPI no rompe la compilación del Back End.

\captura{00-build-success.png}{Compilación exitosa del proyecto.}{fig:build-success}

También se realizó una prueba HTTP autenticada contra el módulo de cuartos. Se creó un cuarto usando credenciales de administrador y posteriormente se consultó la lista de cuartos.

\begin{longtable}{|>{\raggedright\arraybackslash}p{0.30\textwidth}|>{\raggedright\arraybackslash}p{0.30\textwidth}|>{\raggedright\arraybackslash}p{0.25\textwidth}|}
\hline
\rowcolor{cyan!20}
\textbf{Prueba} & \textbf{Endpoint} & \textbf{Resultado}\\
\hline
Crear cuarto & \texttt{POST /api/v1/cuartos} & HTTP 201.\\
\hline
Listar cuartos & \texttt{GET /api/v1/cuartos} & HTTP 200.\\
\hline
\end{longtable}

\captura{07-prueba-http-cuartos.png}{Prueba autenticada del módulo de cuartos.}{fig:prueba-http}

\section{Conclusiones}

La integración de Swagger/OpenAPI permitió documentar automáticamente el Back End del Sistema de Gestión de Cuartos y Reservaciones. La interfaz Swagger UI facilita la revisión de rutas, métodos HTTP, cuerpos de solicitud, respuestas y módulos disponibles sin revisar manualmente todos los controladores.

El proyecto cumple con los puntos solicitados para la práctica: documentación de cuartos, reservaciones, autenticación y seguridad; acceso público a Swagger UI; generación de OpenAPI JSON; y evidencia de compilación y consumo de endpoints protegidos mediante HTTP Basic.

\section{Referencias bibliográficas}
\color{colorESCOM}{
\begin{thebibliography}{9}
\bibitem{spring-boot}
Spring, \textit{Spring Boot Reference Documentation}.
\url{https://docs.spring.io/spring-boot/}

\bibitem{spring-security}
Spring, \textit{Spring Security Reference}.
\url{https://docs.spring.io/spring-security/reference/}

\bibitem{springdoc}
Springdoc, \textit{OpenAPI 3 Library for Spring Boot}.
\url{https://springdoc.org/}

\bibitem{openapi}
OpenAPI Initiative, \textit{OpenAPI Specification}.
\url{https://spec.openapis.org/oas/latest.html}
\end{thebibliography}
}

\end{document}
